\documentclass[12pt, a4paper]{article}

\usepackage[T1]{fontenc}
\usepackage[utf8]{inputenc}
\usepackage{palatino}
\usepackage[protrusion=true,expansion=false]{microtype}
\usepackage[margin=1.3in]{geometry}
\usepackage{setspace}
\usepackage{parskip}
\usepackage{titlesec}
\usepackage{hyperref}
\usepackage{amsmath}
\usepackage{amssymb}
\usepackage{amsthm}
\usepackage{mathtools}
\usepackage{booktabs}
\usepackage[numbers]{natbib}

\hypersetup{
  colorlinks=true,
  linkcolor=black,
  citecolor=black,
  urlcolor=black
}

\titleformat{\section}{\normalfont\large\scshape}{\thesection.}{1em}{}
\titleformat{\subsection}{\normalfont\normalsize\itshape}{\thesubsection.}{1em}{}

\setstretch{1.15}
\DeclarePairedDelimiter{\norm}{\lVert}{\rVert}
\DeclareMathOperator*{\argmax}{arg\,max}

\theoremstyle{plain}
\newtheorem{theorem}{Theorem}[section]
\newtheorem{proposition}{Proposition}[section]
\newtheorem{lemma}{Lemma}[section]
\newtheorem{corollary}{Corollary}[section]

\theoremstyle{definition}
\newtheorem{definition}{Definition}[section]
\newtheorem{remark}{Remark}[section]

\title{\textsc{Preference Fields on Semantic Manifolds}\\
\large Directional Constraints, Spectral Admissibility,\\
and Random Geometric Graph Structure}
\author{Flyxion}
\date{June 2026}

\begin{document}

\maketitle

\begin{abstract}
We develop a geometric theory of preference-guided navigation over corpora
represented as high-dimensional random geometric graphs. A user's evolving
preference is modeled as a time-dependent unit vector $\Phi_t$ in a semantic
embedding space $\mathcal{E} = \mathbb{R}^d$, inducing a scalar preference
field $f_t(v) = \langle \Phi_t, v \rangle$ over the corpus manifold
$\mathcal{M} \subset S^{d-1}$. The admissible neighborhood $A_t(\theta) =
\{i : f_t(v_i) \ge \theta\}$ forms the primary object of study: the set of
corpus items accessible to the current preference state. We establish that
this neighborhood has the structure of a subgraph of a spherical random
geometric graph, with clique probabilities governed by the triangle-penalty
corrections of Hunter, Milojević, and Sudakov. We identify the anisotropic
centroid bias of pretrained embeddings as an inadmissible subspace whose
removal via spectral filtering sharpens the discriminative geometry of the
preference field. The central result is a dimensionless order parameter
$\Lambda_t(\theta) = n_\theta^3 p_\theta^3 / d$ governing a phase
transition: below the critical threshold the admissible neighborhood is
statistically indistinguishable from an Erd\H{o}s--R\'enyi random graph and
geometric structure is undetectable; above it, the spherical geometry is
present and the preference field resolves genuine semantic clusters. This
order parameter unifies the field dynamics, spectral admissibility, and
graph-theoretic clique structure into a single geometric framework applicable
to literature exploration, memory retrieval, recommendation, and semantic
search.
\end{abstract}

\tableofcontents
\newpage

% ─────────────────────────────────────────────────────────────────────────────
\section{Introduction}
% ─────────────────────────────────────────────────────────────────────────────

Many information systems can be understood as the problem of navigating a
high-dimensional semantic manifold under evolving constraints. Classical
approaches emphasize retrieval, ranking, and optimization against a fixed
query. The present work examines a structurally different setting: navigation
driven by a \textit{preference field} that evolves in response to the
navigator's encounters with the manifold itself. The query is not given in
advance; it is constructed through interaction, and the geometry of the space
shapes both what is accessible and what the evolving field can discriminate.

This setting arises whenever users explore a corpus whose internal structure
is represented by embeddings in a high-dimensional metric space. When
preference is encoded as a direction $\Phi_t$ in that space, evaluating
items by $\langle \Phi_t, v_i \rangle$ is not a retrieval operation but a
\textit{field evaluation}: a measurement of how well each point in the corpus
aligns with the current constraint. The natural questions then become
geometric. How does the field evolve? What is the structure of the region it
selects? How does the ambient geometry of the embedding space govern the
possible shapes of coherent preference neighborhoods? When does the geometry
become detectable and meaningful?

We address these questions by developing three interlocking theoretical
structures. The first is a formal framework for preference fields as
directional constraints over manifolds, including a precise account of how
ratings update the field and how the field induces a family of admissible
neighborhoods (Section~\ref{sec:fields}). The second is a random geometric
graph characterization of the high-alignment sub-corpus, drawing on recent
results in Ramsey theory to bound the coherent cluster size and identify the
triangle density as a diagnostic of geometric regime (Sections~\ref{sec:rgg}
and~\ref{sec:threshold}). The third is a spectral admissibility analysis of
the embedding geometry itself, connecting the anisotropy of pretrained
embeddings to a natural notion of inadmissible subspace and characterizing
spectral filtering as admissibility enforcement (Section~\ref{sec:spectral}).

The paper is organized to be progressively more abstract. Early sections
establish concrete objects; later sections develop the connections between
them. We assume throughout that $\mathcal{E} = \mathbb{R}^d$ with $d$ large
(specifically $d = 512$ in the CLIP ViT-B/32 setting, but most results hold
for general $d$) and that items are represented as unit vectors $v_i \in
S^{d-1}$. Section~\ref{sec:discussion} places the framework in the context of
related field theories; no familiarity with those frameworks is required for
the main results.

% ─────────────────────────────────────────────────────────────────────────────
\section{Preference Fields as Directional Constraints}
\label{sec:fields}
% ─────────────────────────────────────────────────────────────────────────────

\subsection{Semantic Embedding and the Corpus Manifold}

Let $\mathcal{E} = \mathbb{R}^d$ be a semantic embedding space equipped with
the standard inner product. A \textit{corpus} $\mathcal{C} = \{c_1, \ldots,
c_N\}$ is a finite collection of items, each associated with a unit vector
$v_i \in S^{d-1}$ via an embedding map $\phi : \mathcal{C} \to S^{d-1}$.

\begin{definition}[Corpus manifold]
The \emph{corpus manifold} $\mathcal{M} \subset S^{d-1}$ is the image
$\phi(\mathcal{C})$. We write $\bar{v} = \frac{1}{N}\sum_{i=1}^N v_i$ for the
corpus centroid (not necessarily of unit norm) and
$\sigma^2 = \frac{1}{N}\sum_{i=1}^N \norm{v_i - \bar{v}}^2$ for the
corpus dispersion.
\end{definition}

The embedding space for pretrained vision-language models (CLIP) and language
models approximates a semantic manifold: inner products $\langle v_i, v_j
\rangle$ correlate with human-perceived similarity between items $c_i$ and
$c_j$~\cite{radford2021clip}. This grounding is what makes the preference
field semantically meaningful; without it the framework reduces to arbitrary
geometry.

\subsection{The Preference Field}

\begin{definition}[Preference field]
A \emph{preference field} at time $t$ is a unit vector $\Phi_t \in S^{d-1}$.
It induces a scalar alignment function
\begin{equation}
  f_t : S^{d-1} \to [-1, 1], \qquad f_t(v) = \langle \Phi_t, v \rangle,
\end{equation}
and its restriction to the corpus manifold
\begin{equation}
  f_t\!\restriction_{\mathcal{M}} : \mathcal{M} \to [-1, 1], \qquad
  f_t(v_i) = \langle \Phi_t, v_i \rangle
\end{equation}
is the \emph{discrete preference field} over $\mathcal{C}$.
\end{definition}

The discrete preference field is not a ranking: it is the trace of a continuous
linear functional on $S^{d-1}$ evaluated at the corpus sample points. This
distinction matters. A ranking selects a permutation of $\mathcal{C}$ with no
structure beyond order. The preference field selects a \textit{direction} in
$\mathcal{E}$, and that direction generalizes continuously to any item whose
embedding can be computed, whether or not it appears in $\mathcal{C}$.

\subsection{Field Dynamics Under Preference Accumulation}

Let $\mathcal{R}_t = \{(i, r_i)\}$ be the set of rated items at time $t$,
where $r_i \in \mathbb{R}$ is a scalar rating with $\mathrm{sgn}(r_i)$
indicating positive or negative preference. Define the weighted sum
\begin{equation}
  \tilde{\Phi}_t = \sum_{(i, r_i) \in \mathcal{R}_t} w(r_i)\, v_i,
  \qquad w(r) = \mathrm{sgn}(r) \cdot (1 + \kappa |r|),
\end{equation}
for a sensitivity parameter $\kappa > 0$. When $\tilde{\Phi}_t \neq 0$, the
updated preference field is
\begin{equation}
  \label{eq:field-update}
  \Phi_t = \frac{\tilde{\Phi}_t}{\norm{\tilde{\Phi}_t}}.
\end{equation}

\begin{proposition}[Monotone convergence under consistent preference]
\label{prop:convergence}
Suppose all ratings in $\mathcal{R}_t$ are positive ($r_i > 0$ for all $i$)
and the rated items lie within a geodesic ball $B_\rho(\Phi^*) \subset
S^{d-1}$ of radius $\rho < \pi/2$ around some direction $\Phi^* \in
S^{d-1}$. Then $\langle \Phi_t, \Phi^* \rangle \ge \cos\rho > 0$ and
$\langle \Phi_t, \Phi^* \rangle$ is non-decreasing as new positively-rated
items are added from $B_\rho(\Phi^*)$.
\end{proposition}

\begin{proof}
The normalized sum of vectors within a geodesic ball of radius $\rho <
\pi/2$ has positive inner product with any vector in the ball. Adding a
new vector $v_j \in B_\rho(\Phi^*)$ to $\tilde{\Phi}_t$ moves the sum
further into the half-space $\{v : \langle v, \Phi^* \rangle > 0\}$,
since $\langle v_j, \Phi^* \rangle > \cos\rho > 0$. Normalization
preserves the sign and the bound follows from the convexity of geodesic
balls on $S^{d-1}$.
\end{proof}

\subsection{Admissible Neighborhoods}

\begin{definition}[Admissible neighborhood]
For threshold $\theta \in [-1, 1]$, the \emph{admissible neighborhood} of
the field $\Phi_t$ at level $\theta$ is
\begin{equation}
  A_t(\theta) = \{i \in \mathcal{C} : f_t(v_i) \ge \theta\}.
\end{equation}
The \emph{inadmissible set} at level $\theta$ is $\mathcal{C} \setminus
A_t(\theta)$.
\end{definition}

As $\theta$ increases, $A_t(\theta)$ shrinks: only items tightly aligned
with the current field direction remain admissible. The trajectory of
$A_t(\theta)$ as $t$ increases characterizes how navigation concentrates the
preference field. In an embedding space with good semantic structure,
$A_t(\theta)$ for large $\theta$ should correspond to a semantically coherent
neighborhood of $\phi^{-1}(\Phi_t)$---the ``concept nearest to the current
field direction.''

\begin{remark}
In practice, $\theta$ is not set explicitly but is induced by a softmax
sampling distribution with temperature $T > 0$:
\begin{equation}
  \mathcal{P}_t(i) \propto \exp\!\left(\frac{s_t(i)}{T}\right),
\end{equation}
where $s_t(i) = \alpha f_t(v_i) + \text{(auxiliary terms)}$ and $\alpha >
0$ bounds the alignment contribution. Items in $A_t(\theta)$ receive
disproportionate sampling probability, inducing a soft version of the
hard thresholding above.
\end{remark}

% ─────────────────────────────────────────────────────────────────────────────
\section{Random Geometric Graph Structure of the Corpus}
\label{sec:rgg}
% ─────────────────────────────────────────────────────────────────────────────

\subsection{The Spherical RGG Model}

We now establish the formal relationship between the preference field
threshold structure and spherical random geometric graphs. The \textit{Gaussian
random geometric graph} $G(n, d, p)$ of Hunter, Milojević, and
Sudakov~\cite{hunter2026gaussian} samples $n$ vectors $x_1, \ldots, x_n
\overset{\text{iid}}{\sim} \mathcal{N}(0, \frac{1}{d}I_d)$ and connects
$x_i, x_j$ if $\langle x_i, x_j \rangle \ge c_p / \sqrt{d}$, where
$c_p$ is chosen so that each edge has marginal probability $p$. This is the
Gaussian analogue of the spherical RGG; in high dimension, Gaussian vectors
concentrate on a sphere of radius $\approx 1$, making the two models
asymptotically close.

\begin{definition}[Threshold graph of the preference field]
For threshold $\theta$, the \emph{threshold graph} $G_t(\theta)$ has
vertex set $A_t(\theta)$ and connects $i, j \in A_t(\theta)$ if
$\langle v_i, v_j \rangle \ge \theta'$ for some proximity threshold
$\theta'$ induced by $\theta$ and the corpus geometry.
\end{definition}

\begin{proposition}[RGG identification]
\label{prop:rgg}
If the corpus embeddings $\{v_i\}$ are drawn from a distribution on
$S^{d-1}$ whose conditional distribution given $\langle v_i, \Phi_t
\rangle \ge \theta$ is approximately uniform on the spherical cap
$\{v \in S^{d-1} : \langle v, \Phi_t \rangle \ge \theta\}$, then
$G_t(\theta)$ is approximately distributed as a spherical RGG on
$|A_t(\theta)|$ vertices with edge probability determined by the
angular diameter of the cap.
\end{proposition}

The ``approximately uniform'' condition holds to a first approximation for
corpora generated by diffusion models conditioned on a fixed prompt, since
the prompt induces a Gaussian-like perturbation around the prompt direction
in CLIP space rather than a hard angular cutoff. The approximation improves
as $d$ grows relative to the angular radius of the cap.

\subsection{Triangle Penalty and Clique Suppression}

The key geometric fact distinguishing an RGG from $G_{n,p}$ is the
\textit{triangle penalty}: edges in the RGG are not independent because
shared proximity to the preference direction correlates pairs of items.
Hunter, Milojević, and Sudakov quantify this precisely for the Gaussian
model.

Let $a = e^{-c_p^2/2}/\sqrt{2\pi}$ where $c_p$ is the inner-product
threshold parameter. Define the per-triangle correction factors
\begin{equation}
  \delta_r = 1 - \frac{a^3 p^3}{\sqrt{d}}, \qquad
  \delta_b = 1 + \frac{a^3(1-p)^3}{\sqrt{d}}.
\end{equation}

\begin{theorem}[Clique probability bounds, Hunter--Milojević--Sudakov~\cite{hunter2026gaussian}]
\label{thm:hms}
Let $G \sim G(n, d, p)$ with $d \ge D^2 \ell^2$ where $D$ is sufficiently
large. Then:
\begin{align}
  P_{\mathrm{red},\ell} &\le p^{\binom{\ell}{2}}
    \exp\!\left(-\frac{a^3 p^3}{\sqrt{d}}\binom{\ell}{3}\right)
    \!\left(1 + O\!\left(\tfrac{1}{D}\right)\right), \label{eq:pred} \\
  P_{\mathrm{blue},\ell} &\le (1-p)^{\binom{\ell}{2}}
    \exp\!\left(\frac{a^3(1-p)^3}{\sqrt{d}}\binom{\ell}{3}\right)
    \!\left(1 + O\!\left(\tfrac{1}{D}\right)\right). \label{eq:pblue}
\end{align}
\end{theorem}

The interpretation in our setting is the following. Color an item $i \in
\mathcal{C}$ ``blue'' if $f_t(v_i) \ge \theta$ (it lies in the admissible
neighborhood) and ``red'' if $f_t(v_i) < \theta$. Blue items are those close
to $\Phi_t$; red items are those far from it. Then:

\begin{itemize}
  \item \textit{Blue cliques} (fully admissible clusters) form with slightly
    higher probability than $G_{n,p}$ predicts ($\delta_b > 1$): proximity
    to the field direction is self-reinforcing, and items aligned with
    $\Phi_t$ tend to be mutually aligned.
  \item \textit{Red cliques} (fully inadmissible clusters) form with
    exponentially lower probability than $G_{n,p}$ predicts: items forced
    far from $\Phi_t$ are correspondingly close to each other, making large
    mutually-distant (inadmissible) sets geometrically unlikely.
\end{itemize}

The exponential suppression in~\eqref{eq:pred} grows as
$\binom{\ell}{3}$---one triangle per triple in the clique---reflecting the
cumulative geometric cost of each triangular constraint. This is the
mathematical expression of the intuition that ``inadmissible regions do not
support large coherent clusters.''

\subsection{Clique Number of the Admissible Neighborhood}

\begin{definition}[Field clique number]
The \emph{field clique number} $\omega_t(\theta)$ is the size of the
largest clique in the threshold graph $G_t(\theta)$: the maximum $k$ such
that $k$ items in $A_t(\theta)$ are mutually within inner-product distance
$\theta'$ of each other.
\end{definition}

\begin{corollary}[Clique bound under geometric constraint]
\label{cor:clique}
Under the RGG identification of Proposition~\ref{prop:rgg}, with
$n_\theta = |A_t(\theta)|$ and edge probability $p_\theta$:
\begin{equation}
  \mathbb{E}[\omega_t(\theta)] \lesssim
  2\log_{1/(1-p_\theta)} n_\theta
  \cdot \left(1 + O\!\left(\frac{1}{\sqrt{d}}\right)\right).
\end{equation}
The correction factor is negative ($\delta_b^{\binom{k}{3}} > 1$ amplifies
rather than suppresses blue cliques), but grows slowly compared to the
independent $G_{n,p}$ upper bound.
\end{corollary}

The corollary establishes that the field clique number is bounded
independently of session length or rating history---it depends only on
the corpus size $n_\theta$ and the edge probability $p_\theta$ at the given
alignment threshold. No amount of additional interaction can cause the
admissible neighborhood to contain a coherent cluster larger than what the
ambient spherical geometry permits. This is a structural, not empirical,
constraint on preference field convergence.

% ─────────────────────────────────────────────────────────────────────────────
\section{Spectral Admissibility of Embedding Geometry}
\label{sec:spectral}
% ─────────────────────────────────────────────────────────────────────────────

\subsection{Anisotropy as Inadmissibility}

Pretrained embedding models produce representations that are
\textit{anisotropic}: the image of the embedding map $\phi : \mathcal{C} \to
S^{d-1}$ is not uniformly distributed over the sphere but concentrates in a
narrow cone~\cite{ethayarajh2019anisotropy}. The centroid of this cone,
\begin{equation}
  \bar{v} = \frac{1}{N}\sum_{i=1}^N v_i,
\end{equation}
is non-zero and can be large in norm relative to the dispersion $\sigma$.
When $\norm{\bar{v}}$ is large, every item has non-negligible inner product
with $\bar{v}$, regardless of semantic content.

This creates a systematic distortion of the preference field. In the
weak-field regime (early in navigation, when $\Phi_t$ has not yet separated
from $\bar{v}$), the alignment $f_t(v_i) = \langle \Phi_t, v_i \rangle
\approx \langle \bar{v}, v_i \rangle$ for most items $i$, making
$f_t\!\restriction_\mathcal{M}$ approximately constant and semantically
uninformative. The field cannot discriminate until $\Phi_t$ has moved
sufficiently far from $\bar{v}$ to produce meaningful variance.

\begin{definition}[Admissible and inadmissible subspaces]
Let $W_U$ be the unembedding matrix of a language or vision-language model,
with singular value decomposition $W_U = U \Sigma V^\top$. The
\emph{edge-spectrum subspace} is
\begin{equation}
  \mathcal{V}_{\mathrm{edge}} = \mathrm{span}\{V_{[k]} : k \in
  \mathcal{K}_{\mathrm{edge}}\},
\end{equation}
where $\mathcal{K}_{\mathrm{edge}}$ indexes the singular dimensions with
either the largest or smallest singular values---the boundary of the
spectrum. The \emph{bulk subspace} is its complement:
$\mathcal{V}_{\mathrm{bulk}} = \mathcal{V}_{\mathrm{edge}}^\perp$.

A vector $v \in \mathbb{R}^d$ is \emph{spectrally admissible} if
$\norm{\Pi_{\mathrm{edge}} v} / \norm{v} < \varepsilon$ for some threshold
$\varepsilon > 0$, where $\Pi_{\mathrm{edge}}$ is the orthogonal projector
onto $\mathcal{V}_{\mathrm{edge}}$.
\end{definition}

Wu et al.~\cite{wu2026embedfilter} establish, via Logit Spectroscopy, that
$\mathcal{V}_{\mathrm{edge}}$ is precisely the subspace responsible for
writing high-frequency, semantically uninformative tokens into the embedding
space. The corpus centroid $\bar{v}$ lies predominantly in
$\mathcal{V}_{\mathrm{edge}}$; removing the edge-spectrum projection
eliminates the centroid bias and redistributes embeddings more isotropically
over $S^{d-1}$.

\subsection{EmbedFilter as Admissibility Enforcement}

\begin{definition}[EmbedFilter transformation~\cite{wu2026embedfilter}]
The \emph{Bulk Spectrum Transformation} with filtering ratio $\tau$ is the
linear map
\begin{equation}
  \Phi_\tau : \mathbb{R}^d \to \mathbb{R}^k, \qquad
  \Phi_\tau(v) = V[l_\tau : r_\tau]^\top v,
\end{equation}
where $k = r_\tau - l_\tau \approx d/\tau$, and $l_\tau, r_\tau$ bound the
retained bulk indices. The \emph{filtered embedding} is
$\tilde{v} = \Phi_\tau(v)$.
\end{definition}

\begin{proposition}[EmbedFilter as inadmissibility removal]
\label{prop:embedfilter}
The EmbedFilter transformation satisfies:
\begin{enumerate}
  \item[(i)] \emph{Projection identity.} $\Phi_\tau = (I -
    \Pi_{\mathrm{edge}}) + \Pi_{\mathrm{bulk}}$, up to
    re-indexing. In particular, $\tilde{v} = (I -
    \Pi_{\mathrm{edge}})v$ up to rotation within the bulk.
  \item[(ii)] \emph{Distance preservation.} For any $v, w \in \mathbb{R}^d$,
    $\norm{\Phi_\tau(v) - \Phi_\tau(w)}_2 = \norm{v - w}_2$ within the
    bulk subspace.
  \item[(iii)] \emph{Centroid elimination.} If $\bar{v} \in
    \mathcal{V}_{\mathrm{edge}}$ then $\Phi_\tau(\bar{v}) = 0$.
  \item[(iv)] \emph{Admissibility enforcement.} Every $\tilde{v}_i =
    \Phi_\tau(v_i)$ is spectrally admissible with respect to the original
    edge-spectrum projector.
\end{enumerate}
\end{proposition}

\begin{proof}
(i) follows from the orthogonality of the singular vectors of $V$.
(ii) follows from the fact that $V[l_\tau : r_\tau]$ consists of columns of an
orthogonal matrix, hence is an isometry on its domain. (iii) holds when the
centroid is exactly in the edge subspace; in practice this is approximate,
and the quality of centroid elimination depends on how concentrated $\bar{v}$
is in $\mathcal{V}_{\mathrm{edge}}$. (iv) is immediate from (i): after
projection, the edge-spectrum component of $\tilde{v}_i$ is zero.
\end{proof}

\subsection{Effect on Preference Field Discriminability}

After applying EmbedFilter, the filtered preference field
\begin{equation}
  \tilde{f}_t(v_i) = \langle \tilde{\Phi}_t, \tilde{v}_i \rangle
\end{equation}
operates in the bulk subspace $\mathcal{V}_{\mathrm{bulk}}$. Because the
centroid bias has been removed, the variance
\begin{equation}
  \mathrm{Var}_i[\tilde{f}_t(v_i)] = \frac{1}{N}\sum_{i=1}^N
  \tilde{f}_t(v_i)^2 - \left(\frac{1}{N}\sum_{i=1}^N \tilde{f}_t(v_i)\right)^2
\end{equation}
is strictly larger than $\mathrm{Var}_i[f_t(v_i)]$ whenever $\Phi_t$ has
a non-negligible edge-spectrum component---that is, whenever the
unfiltered field is being partially driven by the centroid attractor. This
variance increase directly improves the discriminative resolution of the
preference field, allowing meaningful alignment differences to emerge earlier
in the navigation trajectory.

Empirically, Wu et al.\ report up to 14\% improvement in MTEB downstream
task performance across Qwen, Llama, and Mistral backbones at $\tau = 2$
(halving the embedding dimension), confirming that the bulk subspace carries
the semantically relevant structure~\cite{wu2026embedfilter}.

% ─────────────────────────────────────────────────────────────────────────────
\section{The Indistinguishability Threshold and Regime Transitions}
\label{sec:threshold}
% ─────────────────────────────────────────────────────────────────────────────

\subsection{The Signed Triangle Count as a Geometric Statistic}

The central statistic for distinguishing an RGG from $G_{n,p}$ is the
\textit{signed triangle count}. For a graph $H$ on $n$ vertices with
edge-probability $p$ and adjacency matrix $A$, define
\begin{equation}
  T(H) = \sum_{\{i,j,k\} \subset [n]}
  (A_{ij} - p)(A_{jk} - p)(A_{ik} - p).
\end{equation}

\begin{proposition}[Signed triangle count discriminates RGG from $G_{n,p}$]
\label{prop:triangle}
For $G \sim G_{n,p}$ (independent edges): $\mathbb{E}[T(G)] = 0$.
For $G \sim$ spherical RGG with the same marginal edge probability:
\begin{equation}
  \mathbb{E}[T(G)] \approx \frac{n^3 p^3}{\sqrt{d}}.
\end{equation}
In both models $\mathrm{Var}[T(G)] \approx n^3 p^3$.
\end{proposition}

The signal-to-noise ratio for distinguishing the two distributions is
therefore $n^3 p^3 / \sqrt{d} \,/\, \sqrt{n^3 p^3} = \sqrt{n^3 p^3 / d}$.
The geometric signal is detectable when this ratio exceeds a constant, i.e.\
when $d < n^3 p^3$.

\subsection{The Regime Indicator}

\begin{definition}[Preference field regime indicator]
At state $t$, for threshold $\theta$, let $n_\theta = |A_t(\theta)|$ and
let $p_\theta$ be the effective edge probability in $G_t(\theta)$. The
\emph{regime indicator} is
\begin{equation}
  \label{eq:lambda}
  \Lambda_t(\theta) = \frac{n_\theta^3\, p_\theta^3}{d}.
\end{equation}
\end{definition}

$\Lambda_t$ is the natural order parameter of the theory. Every other
quantity---alignment variance, clique number, triangle density, spectral
discriminability---is monotone in $\Lambda_t$. When $\Lambda_t$ crosses its
critical value the system undergoes a phase transition from a regime in which
the corpus geometry is hidden to one in which it is geometrically structured
and navigable. In this respect $\Lambda_t$ plays the same role in preference
field theory that the Reynolds number plays in fluid mechanics, the inverse
temperature $\beta$ plays in statistical mechanics, and the basic reproduction
number $R_0$ plays in epidemiology: it is the single dimensionless parameter
that separates qualitatively distinct phases of the system's behavior.

\begin{proposition}[Field detectability criterion]
\label{prop:detectability}
Let $G_t(\theta)$ be the threshold graph induced by the preference field at
threshold $\theta$. Geometric information about the underlying corpus manifold
is recoverable from $G_t(\theta)$ if and only if $\Lambda_t(\theta)$ exceeds
a constant-order threshold. More precisely:
\begin{itemize}
  \item As $\Lambda_t \to 0$: the distribution of $G_t(\theta)$ converges
    in total variation to $G_{n_\theta, p_\theta}$; no statistical test can
    distinguish the threshold graph from an Erd\H{o}s--R\'enyi random graph;
    the preference field carries no recoverable geometric information.
  \item As $\Lambda_t \to \infty$: the signed triangle count
    $T(G_t(\theta))$ diverges from its $G_{n,p}$ expectation of zero by
    $\Theta(\sqrt{n_\theta^3 p_\theta^3 / d}) \cdot \sqrt{n_\theta^3
    p_\theta^3}$ standard deviations; geometric clustering is detectable
    with power approaching one; the preference field is resolving genuine
    semantic structure.
  \item At $\Lambda_t \asymp 1$: the signed triangle count achieves
    signal-to-noise ratio $\Theta(1)$; this is the critical threshold
    separating detectable from undetectable geometry.
\end{itemize}
\end{proposition}

\begin{proof}[Sketch]
The total variation convergence for $\Lambda_t \to 0$ follows from the
HMS indistinguishability theorem~\cite{hunter2026gaussian}: when $d >
n^3 p^3$ the total variation distance between the spherical RGG and
$G_{n,p}$ goes to zero. The detectability for $\Lambda_t \to \infty$
follows from Proposition~\ref{prop:triangle}: the signal-to-noise ratio
of the signed triangle count is $\sqrt{n_\theta^3 p_\theta^3 / d} =
\sqrt{\Lambda_t}$, which diverges. The critical threshold follows from
setting signal equal to noise.
\end{proof}

\begin{theorem}[Phase transition in preference field geometry]
\label{thm:transition}
Let $\mathcal{C}$ be a corpus with embeddings approximately distributed as a
spherical RGG in dimension $d$. Then:
\begin{itemize}
  \item If $\Lambda_t(\theta) \ll 1$: the sub-corpus $A_t(\theta)$ is
    statistically indistinguishable from $G_{n_\theta, p_\theta}$; the
    preference field cannot resolve geometric structure; the admissible
    neighborhood is a uniformly random subset of size $n_\theta$.
  \item If $\Lambda_t(\theta) \gg 1$: the geometry of $S^{d-1}$ is
    detectable in $A_t(\theta)$; triangles close at a rate exceeding
    $G_{n,p}$ prediction; the admissible neighborhood has genuine
    cluster structure that the preference field is resolving.
\end{itemize}
The transition between regimes occurs at $\Lambda_t(\theta) \asymp 1$,
i.e.\ $n_\theta \asymp (d/p_\theta^3)^{1/3}$.
\end{theorem}

\begin{proof}[Sketch]
The indistinguishability claim for $\Lambda_t \ll 1$ follows from
Milojević, Hunter, and Sudakov's theorem that the total variation distance
between a spherical RGG and $G_{n,p}$ converges to zero when $d > n^3 p^3$
(see~\cite{hunter2026gaussian} and the companion
indistinguishability paper). The detectability claim for $\Lambda_t \gg 1$
follows from Proposition~\ref{prop:triangle}: the signed triangle count
provides a statistically powerful test whenever $\sqrt{n^3 p^3 / d} \gg 1$.
\end{proof}

\subsection{Implications for Navigation}

The regime indicator $\Lambda_t(\theta)$ gives a principled diagnostic of
preference field quality. When $\Lambda_t \ll 1$, the field is not resolving
genuine semantic structure: the admissible neighborhood looks random, and
additional preference signals cannot improve geometric discrimination until
$n_\theta$ grows. When $\Lambda_t \gg 1$, the field is operating in the
geometric regime: the coherent cluster identified by $A_t(\theta)$ is a real
semantic neighborhood, and further preference refinement can sharpen its
boundaries.

This has a practical consequence: the transition $\Lambda_t(\theta) \asymp
1$ marks a natural threshold for switching from exploratory navigation
(increasing $n_\theta$ to find a geometric regime) to refinement
(decreasing $\theta$ to identify the tightest coherent sub-cluster). The
mathematics does not prescribe exploration strategy, but it identifies
where in the parameter space exploration is necessary and where refinement
becomes productive.

% ─────────────────────────────────────────────────────────────────────────────
% ─────────────────────────────────────────────────────────────────────────────
\section{Discussion}
\label{sec:discussion}
% ─────────────────────────────────────────────────────────────────────────────

The central contribution of the preceding sections is that preference-guided
navigation naturally generates a geometric field theory. The preference
vector $\Phi_t$ is not merely a weighted sum of ratings; it is a direction
in $S^{d-1}$ whose induced scalar field has measurable geometric properties.
The admissible neighborhood is not merely a ranking threshold; it is a
subgraph of a random geometric graph whose clique structure and triangle
density encode the semantic organization of the corpus. And spectral
admissibility is not merely an engineering post-processing step; it is the
removal of an inadmissible subspace that systematically distorts the
field's discriminative geometry.

The order parameter $\Lambda_t$ unifies these observations. It is a
single dimensionless quantity computed from the admissible neighborhood
size, the effective edge probability, and the embedding dimension. Every
other observable quantity in the theory is monotone in $\Lambda_t$, and
the phase transition at $\Lambda_t \asymp 1$ separates qualitatively
distinct navigation regimes.

\subsection{Connections to Constraint-Based Field Frameworks}

The mathematical structures developed here bear structural similarity to
several frameworks that study the geometry of accessibility and constraint
in different settings. We record these briefly without claiming formal
equivalence.

The preference field $f_t(v) = \langle \Phi_t, v \rangle$ plays the same
mathematical role as the scalar admissibility field in RSVP-based
formulations~\cite{rsvp2024}: a time-dependent scalar that organizes which
regions of a state space are accessible and shapes the competitive landscape
of local configurations. The spectral admissibility operation $(I -
\Pi_\mathrm{edge})$ acts as a directional entropy-reduction operator
analogous to RSVP's admissibility projection, removing high-entropy
directions that carry no differential semantic information. The present
work neither assumes RSVP nor claims equivalence; the correspondence is at
the level of mathematical role.

The nested sequence of admissible neighborhoods $A_t(\theta_1) \supseteq
A_{t'}(\theta_2) \supseteq \cdots$ as the field sharpens resembles the
recursive, constraint-preserving tiling structure of TARTAN, in which each
refinement step inherits and tightens the structure of the previous level.
Whether the preference field sequence can be expressed as a TARTAN tiling
in a precise categorical sense is an open question.

Preference trajectories $\Phi_1, \ldots, \Phi_T$ induce ordered chains of
admissible projections. Items that persist in $A_t(\theta)$ across many
time steps maintain relevance through continued compatibility with the
evolving field, rather than through static storage---a structure that
resembles Chain of Memory models of persistence~\cite{chainofmemory2024}.
Whether this resemblance can be formalized as a dynamical equivalence or
requires a categorical treatment is also open.

% ─────────────────────────────────────────────────────────────────────────────
\section{Future Directions}
\label{sec:future}
% ─────────────────────────────────────────────────────────────────────────────

\subsection{Adaptive Spectral Admissibility}

The EmbedFilter transformation uses a static admissibility criterion: the
edge-spectrum subspace is identified once from the global unembedding matrix
SVD. A more principled formulation allows the criterion to adapt to the
current field state. Define a parameterized family of admissibility
projectors $\Pi_\theta : \mathbb{R}^d \to \mathbb{R}^k$ and update $\theta$
to maximize the discriminative variance of the filtered field over the rated
set:
\begin{equation}
  \theta_t^* = \argmax_{\theta}\;
  \mathrm{Var}_{i \sim \mathcal{R}_t}\!\left[
  \langle \Pi_\theta(\Phi_t),\, \Pi_\theta(v_i) \rangle \right].
\end{equation}
This is online principal component analysis under the preference-weighted
distribution over $\mathcal{R}_t$, initialized at the static EmbedFilter
solution. The effective embedding space specializes to the semantic cluster
the field currently inhabits, discarding spectral dimensions that are
globally informative but locally irrelevant to the current $\Phi_t$.

\subsection{Continuous Field Theory: From Projector to Connection}

The present framework is fundamentally corpus-relative: the admissibility
projector acts on a finite set of vectors and all objectives are defined
through empirical variance over rated examples. Moving to a continuous
field theory requires the projector to become a geometric object.

The natural formulation replaces the fixed projector $\Pi : \mathbb{R}^d
\to \mathbb{R}^k$ with a smoothly varying family
\begin{equation}
  \Pi_v : T_v\mathcal{M} \to A_v,
\end{equation}
where $A_v \subset T_v\mathcal{M}$ is the admissible tangent subspace at
$v \in \mathcal{M}$. The admissible directions form a distribution $A
\subset T\mathcal{M}$. The preference field is no longer a single vector
$\Phi_t \in S^{d-1}$ but a section
\begin{equation}
  \Phi : \mathcal{M} \to T\mathcal{M}
\end{equation}
whose evolution is constrained to remain inside $A$. The adaptive projector
becomes a \textit{connection} on $A$, specifying how admissible directions
transport as one moves through semantic space:
\begin{equation}
  \frac{D\Pi}{dt} = F(\mathcal{R}_t, \Phi_t),
\end{equation}
where $F$ encodes the local curvature of $\mathcal{M}$.

The fundamental question at this level is whether the distribution $A$ is
\textit{integrable}. By the Frobenius theorem, integrability is equivalent
to the existence of admissible leaves $\mathcal{L} \subset \mathcal{M}$
such that navigation remains within a semantic foliation. An integrable
admissibility distribution implies that semantic concepts correspond to
closed leaves and navigation is path-independent within a leaf. A
non-integrable distribution implies path dependence: the accessible region
depends on the trajectory, not merely the current position. This distinction
separates structurally different classes of semantic corpora: a
highly-organized scientific corpus may possess nearly integrable
admissibility leaves; a heterogeneous internet-scale corpus almost certainly
does not.

\subsection{The Continuous Regime Indicator}

The discrete order parameter $\Lambda_t = n_\theta^3 p_\theta^3 / d$ is
derived from graph statistics over a finite neighborhood. In the continuous
limit we expect it to become a local geometric quantity:
\begin{equation}
  \Lambda(x) = \frac{\rho(x)^3\, p(x)^3}{\kappa(x)},
\end{equation}
where $\rho(x)$ is the local semantic density at $x \in \mathcal{M}$,
$p(x)$ is the local alignment probability, and $\kappa(x)$ is an effective
curvature scale replacing the ambient dimension $d$. The discrete dimension
$d$ measures global indistinguishability; the curvature $\kappa(x)$ measures
local geometric obstruction. The phase transition then becomes not ``enough
vertices versus too many dimensions'' but ``enough local density to overcome
curvature-induced indistinguishability.'' This is a genuine field-theoretic
statement: the transition point is a local property of the manifold, not a
global property of the embedding space.

\subsection{Sheaf-Theoretic Admissibility and Global Sections}

In the continuous formulation, each open neighborhood $U \subset
\mathcal{M}$ possesses its own admissible subspace $A(U) \subset T_U
\mathcal{M}$. These local assignments define a \textit{presheaf of
admissible directions}. The central question becomes whether local
admissibility assignments are consistent: whether sections defined on
overlapping neighborhoods $U_\alpha \cap U_\beta$ agree on the overlap,
and whether local sections can be extended to global ones.

If the admissibility presheaf satisfies the sheaf gluing axioms, then a
global semantic concept exists whenever its local restrictions are
consistent. If it does not---if there are obstructions to gluing---then
the concept is only locally definable and cannot be extended across the
full manifold. The obstruction to gluing is captured by the sheaf
cohomology group $H^1(\mathcal{M}, \mathcal{A})$, where $\mathcal{A}$ is
the sheaf of admissible sections.

Viewed this way, the adaptive admissibility projector is the discrete
precursor of a sheaf-gluing operator. Its discrete version selects
admissible dimensions; its continuous version determines whether local
semantic sections extend into global ones. Concepts whose admissibility
sections glue globally correspond to $H^1 = 0$; concepts with non-trivial
obstruction class correspond to $H^1 \neq 0$ and exist only locally. This
is the natural home of the framework: a geometric theory in which
admissibility is a distribution on a manifold, the projector is a
connection on that distribution, and semantic concepts correspond to
globally gluable sections.

\subsection{Ramsey-Theoretic Stopping Criteria}

Corollary~\ref{cor:clique} gives a Ramsey-theoretic bound on the field
clique number $\omega_t(\theta)$. Once $\omega_t$ stabilizes, the field
has extracted the maximum coherent cluster available at the current
threshold: further preference refinement within the same $\theta$ cannot
produce a larger coherent neighborhood. The Frobenius integrability
condition provides a deeper interpretation: $\omega_t$ stabilization
signals that the field has reached the boundary of the current admissible
leaf. Continued navigation requires either expanding $\theta$ (moving to
a coarser leaf) or generating new corpus points by interpolation in
$\mathcal{E}$ (probing the leaf's interior by synthesis rather than
selection).

\subsection{Formal Functorial Unification}

The structural correspondences of Section~\ref{sec:discussion} suggest a
formal program. Define the \textit{preference field category} $\mathbf{PF}$
whose objects are corpus manifolds equipped with preference fields and whose
morphisms are field-preserving maps. The natural question is whether there
exist functors
\begin{equation}
  F_{\mathrm{RSVP}} : \mathbf{PF} \to \mathbf{RSVP}, \quad
  F_{\mathrm{TARTAN}} : \mathbf{PF} \to \mathbf{TARTAN}, \quad
  F_{\mathrm{CoM}} : \mathbf{PF} \to \mathbf{CoM},
\end{equation}
where $\mathbf{RSVP}$, $\mathbf{TARTAN}$, $\mathbf{CoM}$ are the
categories of the respective frameworks. If such functors exist and are
faithful, the frameworks are provably instances of a single underlying
mathematical structure rather than analogous constructions developed
independently. This program requires careful categorical development of
each target framework and constitutes a natural long-range goal.

% ─────────────────────────────────────────────────────────────────────────────
\section{Mathematical Summary}
\label{sec:summary}
% ─────────────────────────────────────────────────────────────────────────────

For reference, we collect the central definitions and results.

\paragraph{Preference field.} $\Phi_t \in S^{d-1}$; updated by
$\Phi_t = \tilde\Phi_t / \|\tilde\Phi_t\|$ where $\tilde\Phi_t = \sum_{(i,r_i) \in
\mathcal{R}_t} w(r_i) v_i$; induces $f_t(v) = \langle \Phi_t, v \rangle$.

\paragraph{Admissible neighborhood.} $A_t(\theta) = \{i : f_t(v_i) \ge \theta\}$;
converges under consistent preference (Proposition~\ref{prop:convergence}).

\paragraph{Triangle corrections.} $\delta_r = 1 - a^3 p^3 / \sqrt{d}$,
$\delta_b = 1 + a^3(1-p)^3/\sqrt{d}$; clique probabilities in the threshold
graph are governed by Theorem~\ref{thm:hms}.

\paragraph{Spectral admissibility.} $(I - \Pi_{\mathrm{edge}})$ removes the
centroid-biased subspace; EmbedFilter enforces admissibility
(Proposition~\ref{prop:embedfilter}); variance of $\tilde f_t$ over the corpus
increases after filtering.

\paragraph{Regime indicator.} $\Lambda_t(\theta) = n_\theta^3 p_\theta^3 / d$;
$\Lambda_t \ll 1$ is the geometric regime transition
(Theorem~\ref{thm:transition}).

\paragraph{Field clique bound.} $\mathbb{E}[\omega_t(\theta)] \lesssim
2 \log_{1/(1-p_\theta)} n_\theta$ with $O(1/\sqrt{d})$ correction
(Corollary~\ref{cor:clique}).

\bigskip
\noindent\rule{\textwidth}{0.4pt}

\begin{thebibliography}{99}

\bibitem{radford2021clip}
Alec Radford, Jong Wook Kim, Chris Hallacy, Aditya Ramesh, Gabriel Goh,
Sandhini Agarwal, et~al.
\newblock Learning transferable visual models from natural language supervision.
\newblock In \textit{Proceedings of ICML}, 2021.

\bibitem{hunter2026gaussian}
Zach Hunter, Aleksa Milojević, and Benny Sudakov.
\newblock Gaussian random graphs and Ramsey numbers.
\newblock arXiv:2512.17718 [math.CO], 2026.

\bibitem{wu2026embedfilter}
Songhao Wu, Zhongxin Chen, Yuxuan Liu, Heng Cui, Cong Li, and Rui Yan.
\newblock Your unembedding matrix is secretly a feature lens for text
  embeddings.
\newblock arXiv:2606.07502 [cs.CL], 2026.

\bibitem{ethayarajh2019anisotropy}
Kawin Ethayarajh.
\newblock How contextual are contextualized word representations?
\newblock In \textit{Proceedings of EMNLP-IJCNLP}, 2019.

\bibitem{rsvp2024}
Flyxion.
\newblock \textit{RSVP: Scalar-Vector-Entropic Field Cosmology}.
\newblock Monograph, 2024.

\bibitem{chainofmemory2024}
Flyxion.
\newblock \textit{Chain of Memory: Selective Persistence and Iterative
  Refinement}.
\newblock Essay, 2024.

\bibitem{pearl1988}
Judea Pearl.
\newblock \textit{Probabilistic Reasoning in Intelligent Systems}.
\newblock Morgan Kaufmann, 1988.

\bibitem{penrose2003}
Mathew Penrose.
\newblock \textit{Random Geometric Graphs}.
\newblock Oxford University Press, 2003.

\end{thebibliography}

\end{document}
